% 第 15 章　哈密顿力学：在状态空间中航行
\chapter{哈密顿力学：在状态空间中航行}

\step{反常识开场}

\begin{mainline}
先看一张照片。一个荡秋千的孩子悬在半空，秋千绳绷得笔直。现在问一个看似幼稚的问题：下一秒，他是升高还是降低？

你答不出来。请注意，这不是因为你物理学得不够多——就算把他的体重、秋千绳的粗细、当天的天气预报全都拿来，你还是答不出来。因为答案根本不在照片里。秋千停在同一个位置，可以正向左飞，也可以正向右飞，两种未来的开头长得一模一样。

再给你一段半秒钟的录像，问题立刻解决。录像多给了什么？多给了“他正在往哪儿去、多快”。原来，预言未来需要的档案有两栏：第一栏是“在哪儿”，第二栏是“往哪儿去得多快”。少了第二栏，第一栏再精确也没用。

接下来的这件事更加违反常识。想象两只一模一样的台球，此刻在同一张球桌的同一个位置、以同一个速度滚动。第一只球是被一只手匀速推到这个位置的；第二只球是从高处的斜坡上滚下来，恰好此刻滑到这个位置。两只球的历史完全不同——一个“出身”于手掌，一个“出身”于斜坡。问：从今往后，它们的运动会有差别吗？

答案是：不会有任何差别。只要此刻的位置和速度相同，未来就完全相同。大自然不记日记，它只读“现在”这一页。

还要特别指出档案里\emph{没有}的一栏：加速度。你可能会想，要不要把“此刻正在多快地变速”也记下来？不必。力通常由位置决定（弹簧看伸长量，引力看距离），位置给定，加速度就跟着定了。所以两栏就够，第三栏是冗余的。“状态只需位置和动量”这件事，本身就是一个深刻的物理事实，而不是图省事。

生活里其实到处有这个原理的影子。导航软件预计到达时间，只需要你此刻的位置和车速，从不过问你早上几点出门、中途加过几次油。下棋也是：要接着下完一盘棋，只需要当前的局面（再加上轮到谁走），之前的棋谱看一眼少一眼——结局的走向已经写在此刻的盘面上。体检报告和心电监护仪的差别也类似：报告单记下你某个早晨的“位置”，监护仪还盯着各项指标正在“往哪变”；医生判断病情走向，两样都要。物理学家给自然界的这种脾气起了个名字：演化是“无记忆”的——未来只依赖现在的状态，不依赖到达这个状态的路径。

我们从小习惯的世界观是“原因在过去”：考试考砸了，是因为昨晚没复习；感冒了，是因为前天淋了雨。可力学偏偏说：给定此刻的位置与速度，过去的一切可以整本扔掉。这个又傲慢又简洁的主张到底对不对？它凭什么成立？“位置加动量”这份两栏的档案，能不能有一个正式的名字、一座像样的房子？这一章就要给这份档案起名叫“状态”，给它盖一座叫“相空间”的房子，并找出房子里每一条路的交通规则——那就是哈密顿方程。
\end{mainline}

\step{先猜一猜}

照例，先押上你的直觉，再往下看。猜错比猜对更有用。三道题分别对应本章的三根支柱：状态、演化，以及哈密顿量与能量之间的微妙关系。

第一题。两把完全相同的秋千，用完全相同的摆长，在同一时刻恰好都经过最低点。可接下来，一个孩子荡得很高，另一个只荡起一点点。这可能吗？如果可能，差别藏在哪里——藏在此刻，还是藏在过去？

第二题。你想预言一颗小行星一年后的位置。天文学家只肯告诉你它此刻的位置和速度，拒绝提供它过去一百年的轨道记录。你吃亏了吗？

第三题。旋转木马以恒定的角速度转动，木马地板上有一道光滑的直槽，槽里一颗滚珠来回滑动。骑在木马上的人用“相对木马的位置和速度”给滚珠记账。他算出的“动能加势能”守恒吗？如果有什么量守恒，那个量还是不是能量？

写下你的三个答案和理由，本章结尾逐一对账。理由请写得具体些——直觉错在哪一步，往往比答案本身更有教益。

\step{动手实验}

\begin{experiment}[一根钥匙摆的“地图”]
\textbf{材料}：一根约三十厘米的细线、一串钥匙（约五十克）、一支笔、一张纸。有手机的话，用慢动作录像更好。

\textbf{步骤一}：用细线拴住钥匙串，手提线头做成一个摆。把它拉到左侧约三十度角的位置释放，盯住最低点：它冲过最低点时最快，方向向右。

\textbf{步骤二}：让它自由摆动几个回合。等它从右侧荡回来、再次经过最低点时，注意观察：位置和上一次经过时完全一样，但运动方向反了。\emph{同一个位置，装着两种不同的运动。}这就是开场那张照片回答不了问题的原因。

\textbf{步骤三}：画地图。在纸上画一条横轴，标上“偏离角度”（向右偏为正，向左偏为负）；再画一条纵轴，标上“运动的势头”（向右运动为正，向左为负）。摆从左侧释放时：角度为负、势头为零——在图上点一个点。冲过最低点：角度为零、势头最正——再点一个点。到达右端：角度最正、势头为零——第三个点。回到最低点：势头最负——第四个点。把这四个点光滑地连起来，你得到的接近一个闭合的椭圆！

\textbf{步骤四}：换一个小一些的幅度（比如十五度角）释放，再画一圈。你会发现：大圈套小圈，两圈永不相交。

\textbf{步骤五}（选做）：用手机慢动作录像拍下摆动，逐帧记下它冲过最低点的快慢，和你纸上的圈对比。你会发现一个定量规律：幅度减半时，圈在横轴和纵轴方向上都大约缩一半，面积缩到四分之一。圈的面积里藏着某种“守恒的东西”——本章末尾的估算会用到它。

\textbf{记录}：为什么轨迹是闭合的圈？为什么大小不同的圈互不相交？如果你的摆慢慢停下来，圈又会变成什么样？先写下观察，解释在后面。这张“角度—势头图”，就是本章主角“相空间”的毛坯房。
\end{experiment}

还有一个零成本的观察：下次荡秋千或看孩子荡秋千时，留意两个极端时刻——经过最低点的一瞬间人最“冲”，到达两端的一瞬间人最“稳”。运动在“位置最大、势头为零”和“位置为零、势头最大”两种状态之间循环交替。你马上就会看到，这个循环在地图上恰好是一圈闭合的轨迹。这个交替还藏着第 16 章的伏笔：位置与势头像两个互相推让的伙伴，一个最大时另一个恰好为零——这正是“振动”的签名。

\step{旧模型遇到麻烦}

第 6 章给了我们一条功勋卓著的规则：合力决定加速度，也就是位置对时间的“二阶变化率”。用它算过抛体、摆和卫星，全都灵验。可仔细端详它，会发现几处别扭。

麻烦一：\textbf{位置只装了半份档案}。方程写在纸上，想让它开始预言，你必须同时告诉它两件事：初始位置\emph{和}初始速度。只给位置，方程根本无法启动——这正是开场照片的困境。一个方程要两份档案才能开工，说明“位置”这个变量天然只装了状态的一半。

麻烦二：\textbf{速度和位置地位不平等}。在牛顿的语言里，位置是主角，速度只是“位置的变化率”，像个随从。可实验步骤三告诉我们，档案的两栏都是独立的信息：知道位置推不出速度，知道速度也推不出位置。凭什么一个是主、一个是仆？

麻烦三：\textbf{变量一换，熟人就变脸}。第 14 章学会了用角度这样的广义坐标描述摆。可一旦坐标换成角度，“质量乘速度”就不再是那个守恒好用的动量——角度的搭档变成了角动量。在牛顿语言里，哪些量守恒、哪些不守恒，藏在一堆坐标换算的背后，看不清楚。

麻烦四：\textbf{二阶方程把“下一步”藏了起来}。牛顿方程直接告诉你的是加速度——速度的“变化率的变化率”。而实际发生的运动是一小步接一小步挪出来的。我们想要的是这样一条规则：给定现在的完整状态，直接说出下一瞬间状态变成什么。它应该是“一阶”的、位置与动量平起平坐的。

最后还有一本散账。第 9 章和第 10 章分别学会了记两本账：能量守恒一本，动量守恒一本。碰上约束多的系统——滑块叠滑块、双摆、机械臂——两本账外加一堆约束力方程，越翻越乱。能不能把“能量”和“动量”揉进同一个函数，让守恒量自己站到台前、约束力退到幕后？

拉格朗日力学（第 14 章）解决了“选对坐标”的问题，但它的方程仍然是二阶的。还差最后一步：把“往哪儿去”提拔为和“在哪儿”平级的坐标。完成这一步的人，是哈密顿。

\step{发明新概念}

\begin{mainline}
\textbf{第一个新概念：广义动量。}第 14 章给每个系统挑了最合适的坐标，叫广义坐标，记作 \(q\)。现在给每个 \(q\) 配一个搭档，叫广义动量（也叫正则动量），记作 \(p\)。这个搭档是谁，由坐标决定。用直角坐标 \(x\) 描述自由粒子，搭档就是老熟人 \(p=mv\)，质量乘速度。用角度 \(\theta\) 描述摆，搭档是角动量：质量乘摆长的平方、再乘角速度。直觉上，\(p\) 回答的问题不是“跑多快”，而是“想让坐标 \(q\) 增加，得付出多少运动的份量”。花样滑冰运动员收臂时转得更快，守的正是角动量这份“广义动量”：手臂伸开角速度慢，收拢角速度快，两者的乘积几乎不变（第 11 章）。这是“角度的搭档是角动量”在生活中最优雅的一次亮相。

广义动量与普通动量何时相同、何时不同？在直角坐标、没有磁场、没有会动的约束这些“老实”场合，它就是 \(mv\)，和你第 10 章学的动量一模一样。但在三种场合它会变脸：坐标本身是角度时，它是角动量；有磁场时，带电粒子的广义动量要在 \(mv\) 上再加一项磁场的贡献（第 22 章会见面的伏笔）；用了随时间变化的坐标系时，它还要带上坐标系自身的“马达参数”。判据只有一句：\(p\) 永远跟着坐标 \(q\) 走，坐标一变，搭档就换。

\textbf{第二个新概念：相空间。}以 \(q\) 为横轴、\(p\) 为纵轴画一张地图，这张地图叫相空间。地图上的一个点，就是一份完整的状态档案：在哪儿、往哪儿去，两栏俱全。时间流逝，这个点在地图上走出一条轨迹；系统的全部经典命运，就是这条轨迹。

相空间有三条铁律，你在实验里已经见过两条。第一，轨迹永不相交：同一个点只有一条出路，因为同一份状态只有一个未来（这就是步骤四里大小两圈不相交的原因）。第二，轨迹不会凭空产生或消失：点永远连续地流动。第三，每种系统在地图上画出自己特有的图案：摆画闭合的圈，自由粒子画直线，被摩擦拖住的摆画螺旋。将来你还会见到更野性的图案：双摆这样的系统，轨迹可以在有限区域里永远游荡、永不重复，却同样遵守这三条铁律。
\end{mainline}

这个比喻要交代清楚。相空间像什么？像导航软件里的地图：只要标出“你在哪、正往哪开”，整条路线就确定了。它不像什么？它不像真实的地面：地图上的“高处”不是海拔，而是动量；你不能把汽车开进相空间，它只是账本，不是新发现的空间。地图上一个点的“上下移动”，对应的是动量的增减，不是物体真的在升降。

\begin{center}
\begin{tikzpicture}[scale=1.05]
  \draw[-{Stealth}] (-3.7,0) -- (3.7,0) node[right=2pt] {角度 \(q\)};
  \draw[-{Stealth}] (0,-2.0) -- (0,2.0) node[above=2pt] {动量 \(p\)};
  % 外环（大幅度摆动），顺时针流动
  \draw[thick,blue!70!black] (0,0) ellipse [x radius=2.8,y radius=1.4];
  \draw[thick,blue!70!black,-{Stealth}] (0,0) ++(150:2.8 and 1.4) arc[start angle=150,end angle=30,x radius=2.8,y radius=1.4];
  \draw[thick,blue!70!black,-{Stealth}] (0,0) ++(330:2.8 and 1.4) arc[start angle=330,end angle=210,x radius=2.8,y radius=1.4];
  % 内环（小幅度摆动）
  \draw[thick,red!70!black] (0,0) ellipse [x radius=1.6,y radius=0.8];
  \draw[thick,red!70!black,-{Stealth}] (0,0) ++(150:1.6 and 0.8) arc[start angle=150,end angle=30,x radius=1.6,y radius=0.8];
  % 同一位置、两种状态
  \draw[dashed,gray] (1.0,-1.75) -- (1.0,1.75);
  \fill (1.0,1.31) circle[radius=1.7pt];
  \fill (1.0,-1.31) circle[radius=1.7pt];
  \node[right=2pt,align=left] at (1.05,1.31) {A：同一位置，正向右飞};
  \node[right=2pt,align=left] at (1.05,-1.31) {B：同一位置，正向左飞};
  \node[blue!70!black] at (-2.2,1.05) {大幅度};
  \node[red!70!black] at (-1.15,0.55) {小幅度};
\end{tikzpicture}
\end{center}

\begin{mainline}
\textbf{第三个新概念：哈密顿量。}现在给这张地图配上“地形”。把系统的总能量不用速度写、改用 \(q\) 和 \(p\) 写出来，得到的函数叫哈密顿量，记作 \(H(q,p)\)。以摆为例：动能部分写成动量的平方除以两倍的“转动惯量”，势能部分照旧，于是
\[
  H = \frac{p^{2}}{2ml^{2}} + mgl\,(1-\cos\theta)
\]
其中 \(m\) 是质量，\(l\) 是摆长，\(\theta\) 是偏角，\(p\) 是与角度搭档的角动量。它像什么？像地图上的等高线地形图：能量守恒时，轨迹恰好趴在一条等高线上，幅度大的圈等高线高，幅度小的圈等高线低。它不像什么？它不像一座真的山：点沿等高线走，而不是往山下滚——相空间里的流动规则跟爬山完全是两回事。

关键问题：哈密顿量是不是总能量？\emph{在本书大多数场合，是的}——只要约束不随时间变化、势能只依赖位置、你用的坐标系本身不随时间运动，\(H\) 就等于动能加势能。但这不是定律，而是有条件的巧合。反例马上就来：第三题里旋转木马上的人，用随木马转动的坐标记账，会发现滚珠的“动能加势能”忽大忽小，并不守恒；守恒的是另一个组合量——那个组合量正是他这个坐标系里的 \(H\)，它等于总能量减去一项与木马转速有关的修正。他的坐标系自带“马达”，账本上就要扣掉马达的那一份。所以请记住护栏：哈密顿量是一个用 \((q,p)\) 写成的、负责驱动演化的函数；它在很多场合碰巧等于总能量，但不要把“碰巧”当成“永远”。

哈密顿写法还有一个立刻兑现的红利：守恒量变得一目了然。如果某个坐标压根不出现在 \(H\) 里——比如自由粒子的 \(H=p^{2}/(2m)\) 里找不到位置 \(x\)——那么 \(H\) 沿这个方向的坡度为零，对应的动量就一点都不变：动量守恒。反之，摆的角度明明白白写在势能里，角动量就注定不守恒。一句话：\emph{函数里不出现的坐标，配一个守恒的动量}。坐标选择的艺术（第 14 章）和守恒律的发现，在这里握了手；至于这握手背后更深层的原因——对称性——要留到第 58 章揭晓。
\end{mainline}

\step{一句话模型}

\begin{oneline}
把“在哪儿”（\(q\)）和“往哪儿去”（\(p\)）平起平坐地写成一个点，系统的全部经典命运就是这个点在相空间里的流动；哈密顿量 \(H\) 既是这张地图的地形，也是推动点前进的发动机。
\end{oneline}

前半句说的是相空间：状态即点，演化即流动。后半句说的是哈密顿方程的直觉：知道地形（\(H\) 在每个方向上的坡度），就知道点下一步往哪儿流。一座山的地形决定溪水的流向，一个系统的 \(H\) 决定状态的流向——不同的是，相空间里的“溪水”不往低处跑，而是绕着等高线转。背下这句话，你就拿到了日后读懂薛定谔方程（第 44 章）的半张门票。

\step{数学翻译器}

\begin{mainline}
把上面的直觉翻译成规则，只需要两句话。第一句：\(p\) 的变化率，等于 \(H\) 沿 \(q\) 方向的坡度的负值——地形在坐标方向上越“陡”，动量变化越快，这正是“力是势能的坡度”的新衣裳。第二句：\(q\) 的变化率，等于 \(H\) 沿 \(p\) 方向的坡度——动量越大坐标跑得越快，这正是“动量驱动位置”的翻版。两句话位置与动量完全对称，只差一个负号。这个负号不是装饰：正是它让轨迹绕成闭合的圈，而不是一路滑向谷底。

还有一个读图要领。相空间里的流动是沿等高线“切着走”，而不是朝低处冲。原因就在这对称结构里：\(q\) 方向的坡度负责驱动 \(p\)，\(p\) 方向的坡度负责驱动 \(q\)，两个驱动方向互相垂直，合起来的运动恰好与等高线平行。这就是为什么摆的轨迹绕圈而不坍缩——也是为什么能量能不增不减地守在轨迹上。
\end{mainline}

\begin{mathbridge}
把两句话写成方程，就是著名的哈密顿方程：
\[
  \dot{q} = \frac{\partial H}{\partial p}, \qquad
  \dot{p} = -\frac{\partial H}{\partial q}
\]
符号 \(\dot{q}\) 表示 \(q\) 随时间的变化率；\(\partial H/\partial p\) 叫偏导数，意思是“只盯着 \(p\) 看，把别的变量当常数，\(H\) 变化多快”。原来是一个二阶方程，现在是两个一阶方程——代价是变量从一个变成两个，收获是位置与动量平级、状态与演化一目了然。

先做三道例行检查。第一道，令势能为零：\(H = p^{2}/(2m)\)。第二个方程给出 \(\dot{p}=0\)——动量守恒；第一个方程给出 \(\dot{q}=p/m\)——动量等于质量乘速度，定义自洽。合起来：自由粒子在相空间里沿水平直线匀速流动，回到第 6 章的第一定律。第二道，方向反转：把 \(p\) 换成 \(-p\)，第一个方程里的 \(\dot{q}\) 反号——同一位置、反向出发，未来整个倒放，与实验步骤二吻合。第三道，看摆：对上面写的 \(H\) 求坡度，\(\partial H/\partial p = p/(ml^{2})\) 给出“角速度等于角动量除以转动惯量”，而 \(-\partial H/\partial\theta = -mgl\sin\theta\) 给出“角动量的变化率等于回复力矩”——牛顿力学对摆的全部结论原样归来。两套语言不是两套互相竞争的自然规律，而是同一套规律的两种记法。

再来认一种图案：弹簧振子的 \(H=p^{2}/(2m)+kx^{2}/2\)。能量守恒 \(H=E\) 在相空间里正是一个椭圆的方程——横半轴由振幅决定，纵半轴由最大动量决定。你在实验里徒手画的那个圈，就是它的作品。这个椭圆如此重要，第 16 章会花整整一章研究它。

再看能量大起来会发生什么。摆动幅度小时，轨迹是椭圆；幅度大到接近翻顶，圈被拉长成奇怪的卵形；一旦能量足够翻过最高点，轨迹不再是闭合的圈，而变成两条永不回头的波浪线——摆变成了永不停歇的“转圈机器”。同一组方程、两种图案，分界线上的那条轨迹对应“恰好爬到最高点、速度恰好为零”的临界状态。相空间把这些信息全部画在一张图上，这是它比一堆孤立公式强的地方。
\end{mathbridge}

\begin{mainline}
最后做一个带真实数字的估算，它会领我们走到量子力学的门口。拿你的钥匙摆：质量约 \(0.05\) 千克，摆长约 \(0.3\) 米，从三十度角释放。势能差 \(E = mgl(1-\cos\theta)\) 大约是 \(0.05\times9.8\times0.3\times0.134\approx0.02\) 焦耳——这就是相空间里那个圈所在等高线的高度。顺手验证一下合理性：这些能量在最低点全部变成动能，\(\frac{1}{2}mv^{2}\approx0.02\) 焦耳，解出 \(v\approx0.9\) 米每秒——和实验里“冲过最低点”的快慢一致。摆的频率约 \(f=\frac{1}{2\pi}\sqrt{g/l}\approx0.9\) 赫兹。用能量除以频率，得到这个圈的“面积”约为
\[
  J = \frac{E}{f} \approx \frac{0.02}{0.9} \approx 0.022 \ \text{焦耳·秒}
\]
这个“相空间环面积”有个专门的量纲：能量乘时间。现在请出量子力学的基本常量——普朗克常量 \(h\approx6.6\times10^{-34}\) 焦耳·秒，量纲一模一样。除一下：\(J/h\approx3\times10^{31}\)。意思是：你这根不起眼的钥匙摆，在相空间里画的那个小圈，面积大得能装下三亿亿亿亿个“量子格”。量子力学将告诉我们，相空间的面积不能无限细分，最小一格大约就是 \(h\)——经典的点，在量子世界里是一块面积为 \(h\) 的模糊斑。日常生活之所以“经典”，正是因为我们的圈太大，大到格子小得看不见。
\end{mainline}

\begin{challenge}
哈密顿量从哪里来？第 14 章有拉格朗日量 \(L(q,\dot{q})\)。定义广义动量 \(p=\partial L/\partial\dot{q}\)，再做一次“换自变量”的手术（数学上叫勒让德变换）：\(H(q,p)=p\dot{q}-L\)，把对 \(\dot{q}\) 的依赖换成对 \(p\) 的依赖，哈密顿方程就自动跳出来。这个手术解释了为什么 \(p\) 有时不是 \(mv\)：有磁场时 \(L\) 里含有速度与磁场耦合的项，求出来的 \(p\) 自然多出一截。

还有两个通往深处的路标。其一，刘维尔定理：相空间里一小团邻近的点，流动时形状会扭曲，但体积严格不变——像一滴不可压缩的墨水在河里被拉伸、打卷，却永不涨缩。这是第 31 章统计物理“数状态”的合法执照，也是加速器束流物理的根基。其二，量子的入口：\(q\) 与 \(p\) 这对搭档有一个深层关系（泊松括号 \(\{q,p\}=1\)）；第 46 章将看到，把括号换成量子对易关系、把 \(H\) 换成哈密顿算符，经典力学就平滑地过渡成量子力学——同一个 \(H\)，从地图的地形变成推动量子态随时间演化的发动机（第 44 章）。
\end{challenge}

\step{模型的边界}

\begin{mainline}
哈密顿力学强大，但请先记住第一条边界：\textbf{它没有给世界增加任何新的力或新的定律}。对本书讨论过的力学问题，它与牛顿力学、拉格朗日力学严格等价——同一套物理，三种记法。它的价值不在“算出新结果”，而在“把结构看清楚”：状态是什么、演化是什么、守恒量在哪里。看清楚的回报要后面才兑现：统计物理、量子力学都直接用这种语言书写。

第二条边界：\textbf{它要求力能写成势能的坡度}。摩擦、空气阻力这样的耗散力做不到这一点。耗散系统在相空间里是什么样子？你的钥匙摆其实已经在演示了：实验中那圈并不是严格的闭合环，而是一圈比一圈小的螺旋，最终缩向原点。相空间图一画，耗散无处遁形——哈密顿的框架装不下它，必须扩展（把空气、桌面这些“环境”也算进系统，或者干脆换统计物理的语言）。反过来说，\emph{轨迹是否收缩，就是检验系统有无耗散的试金石}。再补一句公道话：一杯热咖啡在桌上凉下来，若把“咖啡加空气加桌子”看成一个更大的系统，总账仍然守恒——耗散只是“记账范围划得太小”时看到的表象。第 29 章会把这句话展开。

第三条边界：\textbf{不要把 \(H\) 永远等同于总能量}。条件前面交代过：约束不随时间变化、坐标系本身不带“马达”、势能只依赖位置。旋转木马的例子就是反例；磁场中的带电粒子是另一个——那时 \(p\neq mv\)，于是 \(p^{2}/(2m)\) 不再是动能，\(H\) 的形式与“动能加势能”貌合神离。判据只有一句：检查你选的坐标和约束本身是否随时间变化。

第四条边界：\textbf{相空间的点是经典理想化}。它假设 \(q\) 和 \(p\) 可以同时知道得任意精确。第 43 章将看到，自然不允许：位置和动量的精确程度互相牵制，这种牵制不是仪器粗糙造成的，而是量子态本身的性质。所以相空间里的“点”是一块最小面积约 \(h\) 的斑，相空间这幅清晰的地图，是那块斑太小、看不清时的近似。

最后附一份本章的典型误用清单。把广义动量不分场合地当成 \(mv\)：在有磁场或旋转坐标里会算错守恒量。把相空间当成某种真实空间，追问“相空间在哪里”：它在纸面上和计算机里，不在天上。一看到符号 \(H\) 就代入能量公式：先检查成立条件再动手。以为哈密顿力学能算出牛顿力学算不出的新现象：算不出，它只是把同一本账记得更整齐。
\end{mainline}

\begin{center}
\begin{tikzpicture}[scale=1.0]
  \draw[-{Stealth}] (-3.3,0) -- (3.3,0) node[right=2pt] {\(q\)};
  \draw[-{Stealth}] (0,-2.2) -- (0,2.2) node[above=2pt] {\(p\)};
  \draw[thick,blue!70!black,-{Stealth}] plot[domain=0:14,variable=\t,samples=150]
    ({(1.9-0.125*\t)*cos(\t r)},{(1.9-0.125*\t)*sin(\t r)});
  \node[align=left,gray!50!black] at (2.3,-1.7) {有摩擦的摆：\\轨迹旋向原点};
\end{tikzpicture}
\end{center}

\step{回头解释开场现象}

现在回到开场那张照片。它为什么回答不了“下一秒升高还是降低”？因为照片只提供了档案的第一栏——位置。在相空间地图上看，只给定位置相当于只画了一条竖线：竖线与无数条轨迹相交，每个交点都是一个可能的未来。A 点和 B 点在同一条竖线上，位置完全相同，却一个向上、一个向下，各走各的圈。补上第二栏——速度与方向——竖线缩成一个点，点唯一确定一条轨迹，未来就定了（在经典力学的范围内“定了”）。

第一题的答案也清楚了：两把秋千同一时刻都在最低点，位置一样，但动量不同——地图上是两个不同的点，躺在两个不同大小的圈上，所以一个荡得高、一个荡得低。差别不在过去，就在此刻的状态里。

第二题：不吃亏。小行星的全部轨道历史已经压缩进“此刻的位置和动量”这两个数里，过去一百年的记录除了帮你核对，再无新信息。这就是“自然不记日记”的精确含义。

开场那两只“出身”不同的台球，现在也各就各位：在相空间里，它们此刻是同一个点。历史的两本书都合上了，只剩一条共同的轨迹。

第三题最微妙：骑在旋转木马上的人发现滚珠的能量并不守恒，但有一个守恒的组合量。那个量就是他这套转动坐标里的哈密顿量——守恒，却不等于动能加势能。这个例子值得我们回头再看一眼：\emph{守恒的是什么、能量又是什么，竟然取决于你站在哪种坐标里记账}。这不是说物理规律随人而异，而是说“能量”这个词的含义比你以为的更讲究——第 58 章将揭晓，能量守恒来自“物理规律不随时间变化”这条对称性，届时一切各就各位。

\step{脑洞实验与挑战题}

\begin{thought}[把太阳系装进一个点]
太阳加八大行星：每个行星需要三个位置坐标、三个动量坐标，九个天体一共五十四个数。于是，整个太阳系的状态是五十四维相空间里的\emph{一个点}。万有引力定律给定后，这个点走出一条唯一确定的轨迹——行星的所有会合、冲日、掩食，早已“写”在这条轨迹里。

再放肆一点：把可见宇宙的每个粒子都登记上，大约需要 \(10^{80}\) 个坐标。宇宙的经典状态，就是一个维度高到无法想象的相空间里的一个点。这个画面有个著名的历史版本：拉普拉斯在 1814 年设想，一个知道所有粒子位置与动量、又足够强大的智者，能把宇宙的过去与未来尽收眼底——后人把他叫做“拉普拉斯妖”。它曾是机械决定论的图腾；而相空间的视角恰好让我们看清它的两处软肋。其一，真实宇宙是量子的，“一个点”只是经典近似；其二，即使经典画面成立，像三体问题这样的系统对初始点极其敏感——起点差一根头发丝，轨迹很快分道扬镳，预言的实际期限远没有哲学图景那样乐观。确定性写在相空间里，可预测性却是另一回事。
\end{thought}

这套语言早已走出书房，两个去处值得一提。其一是\textbf{粒子加速器}：上亿个粒子的束流，被当作相空间里一团会流动的点云来管理。工程师布置磁铁的“航道”，让这团点云转几千万圈也不散开；刘维尔定理说点云在相空间里的体积守恒，想把束流“压实”就得把涨落搬走——随机冷却技术正是靠这一招拿下了 1984 年的诺贝尔物理学奖。其二是\textbf{量子力学本身}：第 44 章的薛定谔方程，把经典的 \(H(q,p)\) 换成哈密顿算符，让它去推动量子态随时间演化。本章的地图与发动机，到那里原封不动地继续服役——这正是一百年前物理学家选择哈密顿语言作为通往量子世界跳板的原因。

还有一个去处藏在实验室和超算中心里：数值模拟。天体力学家计算水星轨道几万年的进动，材料学家模拟蛋白质怎样折叠，工程师预测卫星编队的长期演化——计算机里跑的，正是把哈密顿方程一小步一小步离散化的算法：相空间里点的流动，被切成一帧一帧的电影。好的算法会刻意尊重相空间的铁律（轨迹不相交、点云体积守恒），否则算到后面，行星会莫名其妙地飞丢。

\begin{puzzles}
\item 自由粒子在相空间里沿水平直线流动。那么竖直上抛的小球（只受重力），在“高度—动量”相空间里画出什么曲线？上升段和下降段是什么关系？\\
\emph{思路提示}：能量守恒给出 \(p^{2}/(2m)+mgx=E\)，对每个高度 \(x\) 解出两个动量，一正一负。曲线是一条横放的抛物线：上半支是上升，下半支是下降，同一轨迹的两个行进方向。
\item 有人宣称：“摩擦只是错觉，只要把坐标取得足够聪明，阻尼摆也能写进哈密顿方程。”请用相空间图反驳他。\\
\emph{思路提示}：哈密顿流的轨迹互不相交、点云体积守恒；而阻尼摆的所有轨迹都旋向原点，任意一团点云的面积最终缩到零。坐标变换不能把“收缩”变没——收缩是物理，不是视角。
\item 为什么相空间里的轨迹绝对不能相交？如果有一天你在实验数据里发现两条轨迹相交了，应该怀疑什么？\\
\emph{思路提示}：交点意味着同一份状态有两个不同的未来，违背“状态唯一决定演化”。若真的相交，说明你的档案漏了一栏——有某个重要变量（比如粒子的电荷、自旋，或环境的影响）没被记进状态里。
\item 旋转木马上的观测者发现：滚珠的“能量”不守恒，但某个组合量守恒。这个组合量随木马的转速改变。由此说说：“什么是能量”和“什么守恒”这两件事的关系，比中学课本告诉我们的更微妙在哪里？\\
\emph{思路提示}：守恒量由你描述世界所用的坐标是否随时间变化来决定；守恒的是哈密顿量 \(H\)，而 \(H\) 与“动能加势能”的关系取决于坐标的选法。第 58 章会把守恒律与对称性正式连接起来。
\end{puzzles}
